\begin{figure}[htbp]
			\centering
			\resizebox{\textwidth}{!}{%
				\begin{tikzpicture}[
					scale=0.95,
					>=latex,
					trave/.style={line width=1.2pt},
					carico/.style={black, line width=0.4pt},
					forza/.style={black, line width=1pt},
					risultante/.style={gray!65, line width=1.8pt},
					diagT/.style={black, line width=0.8pt},
					diagM/.style={black, line width=0.8pt},
					diagTb/.style={gray!65, line width=1.0pt},
					diagMb/.style={gray!65, line width=1.0pt},
					hatch/.style={gray!40, line width=0.3pt},
					quota/.style={black, line width=0.4pt},
					fibre/.style={gray!70, line width=0.3pt},
					segno/.style={circle, draw=black, line width=0.4pt,
						minimum size=10pt, inner sep=0pt, font=\small\bfseries},
					]
					
					% ── Parametri dei vincoli (convenzione C06/C07) ──
					\def\tw{0.15}    \def\thR{0.28}   \def\rn{0.06}
					\def\rw{0.05}    \def\bw{0.09}    \def\gw{0.22}
					\def\hh{0.10}
					
					% ── Macro vincoli (riuso nei tre pannelli) ──
					\def\cerniera#1{%
						\draw[thick] (#1,0) -- (#1-\tw,-\thR) -- (#1+\tw,-\thR) -- cycle;
						\draw[thick] (#1-\gw,-\thR) -- (#1+\gw,-\thR);
						\foreach \xh in {-0.16,-0.08,0,0.08,0.16}{%
							\draw[thin] (#1+\xh,-\thR) -- ++(-0.07,-\hh);}
						\draw[thick, fill=white] (#1,0) circle (\rn);}
					\def\carrello#1{%
						\draw[thick] (#1,0) -- (#1-\tw,-\thR) -- (#1+\tw,-\thR) -- cycle;
						\draw[thick] (#1-\bw,-\thR-\rw-0.01) circle (\rw);
						\draw[thick] (#1+\bw,-\thR-\rw-0.01) circle (\rw);
						\draw[thick] (#1-\gw,-\thR-2*\rw-0.04) -- (#1+\gw,-\thR-2*\rw-0.04);
						\foreach \xh in {-0.16,-0.08,0,0.08,0.16}{%
							\draw[thin] (#1+\xh,-\thR-2*\rw-0.04) -- ++(-0.07,-\hh);}
						\draw[thick, fill=white] (#1,0) circle (\rn);}
					
					% ── Geometria del problema ──
					\def\Lell{4}
					\pgfmathsetmacro{\xA}{0}
					\pgfmathsetmacro{\xMid}{0.5*\Lell}      % mezzeria ℓ/2 = 2
					\def\yqi{-0.80}
					\def\ylab{-3.70}
					
					% ── Centri utili per le etichette (pre-calcolati per evitare ambiguità TikZ) ──
					\pgfmathsetmacro{\xMidAC}{0.5*\xMid}            % centro tratto AC
					\pgfmathsetmacro{\xMidCB}{0.5*(\xMid+\Lell)}    % centro tratto CB
					\pgfmathsetmacro{\xLabel}{0.5*\Lell}            % x dell'etichetta di pannello
					\pgfmathsetmacro{\yLabel}{0.4*\ylab}            % y dell'etichetta di pannello
					
					% ── Parametri del carico distribuito ──
					\def\pgap{0.10}
					\def\php{0.40}
					\pgfmathsetmacro{\pBase}{\pgap}
					\pgfmathsetmacro{\pTop}{\pgap+\php}
					
					% =========================================================
					% PANNELLO (a) — Schema statico: distribuito + risultante
					% =========================================================
					\begin{scope}[xshift=0cm]
						
						\draw[trave] (0,0) -- (\Lell,0);
						\draw[fibre, dashed] (0, -0.1) -- (\Lell, -0.1);
						
						% Cerniera in A
						\draw[thick] (0,0) -- (-\tw,-\thR) -- (\tw,-\thR) -- cycle;
						\draw[thick] (-\gw,-\thR) -- (\gw,-\thR);
						\foreach \x in {-0.16,-0.08,0,0.08,0.16}
						\draw[thin] (\x,-\thR) -- ++(-0.07,-\hh);
						
						% Carrello verticale in B
						\draw[thick] (\Lell,0) -- (\Lell-\tw,-\thR) -- (\Lell+\tw,-\thR) -- cycle;
						\draw[thick] (\Lell-\bw,-\thR-\rw-0.01) circle (\rw);
						\draw[thick] (\Lell+\bw,-\thR-\rw-0.01) circle (\rw);
						\draw[thick] (\Lell-\gw,-\thR-2*\rw-0.04) -- (\Lell+\gw,-\thR-2*\rw-0.04);
						\foreach \x in {-0.16,-0.08,0,0.08,0.16}
						\draw[thin] (\Lell+\x,-\thR-2*\rw-0.04) -- ++(-0.07,-\hh);
						
						\draw[thick, fill=white] (0,0) circle (\rn);
						\draw[thick, fill=white] (\Lell,0) circle (\rn);
						
						% Riquadro del carico distribuito p̄
						\draw[carico] (0,\pBase) rectangle (\Lell,\pTop);
						\foreach \x in {0.25, 0.75, 1.25, 1.75, 2.25, 2.75, 3.25, 3.75} {
							\draw[carico, -{Latex[length=2pt, width=1.5pt]}]
							(\x, \pTop) -- (\x, \pBase);
						}
						\node[above, font=\small] at (0.5, \pTop+0.02) {$\bar{p}$};
						
						
						\node[above left=-2pt,  font=\small] at (0,0)      {$A$};
						\node[above right=-2pt, font=\small] at (\Lell,0)  {$B$};
						%\node[below=2pt, font=\small] at (\xMid, 0) {$C$};
						
						
						\node[font=\small] at (\xLabel, \yLabel+0.25) {(a)};
						
					\end{scope}
					
					% =========================================================
					% PANNELLO (b) — Diagramma del taglio T(z)
					% =========================================================
					\begin{scope}[xshift=5.5cm]
						
						\draw[trave] (0,0) -- (\Lell,0);
						
						\cerniera{0}
						\carrello{\Lell}
						\node[above left=-2pt,  font=\small] at (0,0)      {$A$};
						\node[above right=-2pt, font=\small] at (\Lell,0)  {$B$};
						
						\def\sT{0.4}
						\pgfmathsetmacro{\Tmax}{(\Lell/2)*\sT}
						
						
						% Hatching del caso (a)
						\foreach \x in {0.10, 0.30, ..., 3.90} {
							\pgfmathsetmacro{\h}{(\Lell/2 - \x)*\sT}
							\draw[hatch] (\x, 0) -- (\x, \h);
						}
						
						% Caso (a) NERO: triangolo lineare
						\draw[diagT]
						(0, 0) -- (0, \Tmax)
						-- (\Lell, -\Tmax) -- (\Lell, 0);
						
						% ── Simboli ⊕/⊖ cerchiati nelle due zone ──
						\node[segno] at (\xMidAC/2, {0.4*\Tmax}) {$+$};
						\node[segno] at (1.20*\xMidCB, {-0.4*\Tmax}) {$-$};
						
						\node[left, font=\small] at (0., \Tmax-0.02) {$\tfrac{\bar{p}\ell}{2}$};
						\node[right, font=\small] at (\Lell, -\Tmax+0.02) {$\tfrac{\bar{p}\ell}{2}$};
						
						\node[font=\small] at (\xLabel, \yLabel+0.25) {(b)};
						
					\end{scope}
					
					% =========================================================
					% PANNELLO (c) — Diagramma del momento M(z)
					% =========================================================
					\begin{scope}[xshift=11.0cm]
						
						\draw[trave] (0,0) -- (\Lell,0);
						\cerniera{0}
						\carrello{\Lell}
						\node[above left=-2pt,  font=\small] at (0,0)      {$A$};
						\node[above right=-2pt, font=\small] at (\Lell,0)  {$B$};
						
						\def\sM{0.38}
						\pgfmathsetmacro{\Mb}{-(\Lell*\Lell/4)*\sM}
						\pgfmathsetmacro{\Ma}{-(\Lell*\Lell/8)*\sM}
						
						
						% Hatching del caso (a) — parabola
						\foreach \x in {0.10, 0.30, ..., 3.90} {
							\pgfmathsetmacro{\h}{(\x*(\Lell-\x)/(\Lell*\Lell/4))*\Ma}
							\draw[hatch] (\x, 0) -- (\x, \h);
						}
						
						% Caso (a) NERO: parabola
						\draw[diagM]
						(0.00, 0)
						\foreach \x in {0.20, 0.40, ..., 3.80} {
							-- ({\x}, {(\x*(\Lell-\x)/(\Lell*\Lell/4))*\Ma})
						}
						-- (\Lell, 0);
						
						\node[above, font=\small] at (\xMid, \Ma+0.02) {$\tfrac{\bar{p}\ell^2}{8}$};
						%\node[above, font=\small, gray!50!black] at (\xMid, \Mb+0.02) 	{$\tfrac{\bar{p}\ell^2}{4}$};
						
						% Tangente al vertice della parabola — parallela alla corda A–B (l'asse).
						% Visualizza la proprietà di tangenza in mezzeria nel caso degenere
						% M(A)=M(B)=0: corda coincide con l'asse, tangente in mezzeria orizzontale.
						\draw[diagMb, dashed] (\Lell/5,\Ma) -- (4*\Lell/5,\Ma);
						
						\node[font=\small] at (\xLabel, \yLabel+0.25) {(c)};
						
					\end{scope}
					
				\end{tikzpicture}
			}
			\caption{Problema ``0'': carico uniforme
				sul sistema principale (a); taglio
				$T_0$ (b); momento $M_0$ (c).}
			\label{fig:stat_iper_trasv_01-p0}
		\end{figure}
